\documentclass{UnBeamer}%
%\documentclass[draft]{UnBeamer}%
%\documentclass[handout]{UnBeamer}%
% \documentclass[notasDeAula]{UnBeamer}%

\usepackage{../ED}%
\usepackage{listings}% Código
\subtitle{Apresentação}%

\begin{document}%
	\frame[label=main]{\titlepage}%

	\section{Apresentação}
	\subsection{Objetivos do Curso}
	\begin{frame}[label=main]
	A disciplina 116319 - Estruturas de Dados busca propiciar aos alunos conceito sólidos na forma de  armazenamento e organização dos dados por meio do gerenciamento dinâmico da memória e de estruturas de dados lineares e não lineares. Com a compreensão desses conceitos é esperado o desenvolvimento do pensamento computacional e das habilidades de programação pelos alunos através das atividades práticas.

		\Large Apresentar técnicas fundamentais na análise e projeto de estruturas
		de dados em programas e computacionais.
		\begin{center}
			\includegraphics[width=.6\textwidth]{hammer-screw}
		\end{center}
	\end{frame}

	\subsection{Obrigatória x Optativa}
	\begin{frame}[label=main]
		\begin{center}
			\includegraphics[width=.75\textwidth]{Theofilos_monomaxia_Achilleos}\\
			Pintura de Theophilos Hatzimihail
		\end{center}
	\end{frame}

	\subsection{Aulas Expositivas}
	\begin{frame}[label=main]
		\vspace{-.5\baselineskip}
		\begin{center}
			\includegraphics[width=.75\textwidth]{mafalda}
		\end{center}
		Quino: \url{http://www.quino.com.ar}
	\end{frame}

	\subsection{Na prática, a teoria é outra... }
	\begin{frame}[label=main]
		\begin{center}
			\includegraphics[width=\textwidth]{matriculaweb}
		\end{center}
	\end{frame}

	% \subsection{A1-10/28}
	% \begin{frame}[label=main]
	%	\begin{columns}
	%		\column{.5\textwidth}\centering
	%			
\includegraphics[height=8\baselineskip]{thumbup}
	%		\column{.5\textwidth}\centering
	%			
\includegraphics[height=8\baselineskip]{forbid}
	%	\end{columns}
	% \end{frame}

	\subsection{Ferramentas}
	\begin{frame}[fragile,label=main]
		\begin{columns}[T]
			\column{.75\textwidth}
				\begin{itemize}
					\item Editor de texto
					\item gcc - \emph{GNU Compiler Collection}
					\item Moodle\footnote[frame]{\url{aprender.unb.br/course/view.php?id=2424}}
						\subitem{senha: \alert{ED-1-2013}}
				\end{itemize}
			\column{.2\textwidth}
				\vspace{-1\baselineskip}
				
\includegraphics[width=\textwidth]{gcc}%
		\end{columns}
		\vfill
		\begin{columns}
			\column{.43\textwidth}
				\begin{lstlisting}[language=C,title=hello.c,morekeywords={printf}]
					#include <stdio.h>

					int main() {
					  printf("\nHello World!\n\n");
					  return 1;
					}
				\end{lstlisting}
			\column{.5\textwidth}
				\includegraphics[width=\textwidth]{hello}%
		\end{columns}
	\end{frame}

	\subsection{Definições}
	\begin{frame}[label=main]
		\begin{exampleblock}{Algoritmo}
			\pause
			Sequência \alert{finita} de \emph{instruções} para executar uma tarefa.
			\begin{itemize}
				\item Bem definidas e não ambíguas
				\item Executáveis com uma quantidade de esforço finita
				\item Executáveis em um período de tempo finito
			\end{itemize}
		\end{exampleblock}
		\vfill\pause
		\begin{block}{Estrutura de Dados}
			\pause
			Modo particular de armazenar/organizar dados para facilitar acesso e modificações.
			\subitem{\alert<4>{Só tem significado quando \emph{associada a um algoritmo}.}}
		\end{block}
		\pause
		\begin{columns}[T]
			\column{.68\textwidth}
				Representação dos dados depende de:
				\begin{itemize}
					\item tipo de dado
					\item operações a serem realizadas sobre os dados
					\item etc.
				\end{itemize}
			\column{.3\textwidth}\centering
				% \uncover<3->{
				% \begin{tikzpicture}[node distance=1em,minimum width=5.5em, minimum height=1.5em]
				%	\node[block] (percepcao) {\scriptsize Percepção};
				%	\node[block, below=of percepcao] (processamento) {\scriptsize Processamento};
				%	\node[block, below=of processamento] (execucao) {\scriptsize Execução};
				%	\draw[line, ->] (percepcao) -- (processamento);
				%	\draw[line, ->] (processamento) -- (execucao);
				%	\draw[line, ->] (execucao)..controls +(east:2) and +(east:2) ..  (percepcao);
				% \end{tikzpicture}
				% }
		\end{columns}
	\end{frame}

	\begin{frame}[label=main]
		\frametitle{Estruturas de Dados}
		\begin{center}
			\includegraphics[width=.75\textwidth]{swiss_army}\\
			\href{http://secure.victorinox.com/}{Victorinox Presentation Master}
			\end{center}
	\end{frame}
\end{document}%
